% © 2026 D. Jaroš — CC BY-NC-ND 4.0
%
% File: research_tracks/EW/lepton_mass_ratios_ubt.tex
% Purpose: Documents structural coincidences between lepton mass ratios
%          and UBT constants alpha, Omega_eta, C_Q.
%
% Key findings:
%   h_tau ≈ sqrt(2)*alpha  (-1.0%)  [MC+]
%   m_tau/m_e ≈ 2pi/Omega_eta  (-3.8%)  [MC+]
%   m_tau/m_mu ≈ 4*alpha/Omega_eta  (-7.5%)  [MC+]
%
% Status: [MC+] structural evidence. Not yet derived from S[Theta].
%         VEV hierarchy (M_EW << M_UBT) remains the blocking open problem.
% Date: 2026-06-14

\ifdefined\INCLUDEMODE\else
\documentclass[12pt,a4paper]{article}
\usepackage[a4paper,margin=2.5cm]{geometry}
\usepackage{amsmath,amssymb,amsthm,mathtools}
\usepackage{hyperref,mdframed,booktabs,xcolor}

\newtheorem{theorem}{Theorem}[section]
\newtheorem{proposition}[theorem]{Proposition}
\theoremstyle{remark}\newtheorem{remark}[theorem]{Remark}
\theoremstyle{definition}\newtheorem{observation}[theorem]{Observation}

\newmdenv[backgroundcolor=green!12,linecolor=green!60!black,linewidth=1.5pt]{resultbox}
\newmdenv[backgroundcolor=yellow!15,linecolor=orange,linewidth=1.5pt]{mcbox}
\newmdenv[backgroundcolor=red!8,linecolor=red!60!black,linewidth=1.5pt]{openprobbox}

\newcommand{\Lz}{\textbf{[L0]}}
\newcommand{\Lo}{\textbf{[L1]}}
\newcommand{\MC}{\textbf{[MC+]}}
\newcommand{\OPEN}{\textbf{[Open]}}

\title{\textbf{Lepton Mass Ratios from UBT Constants}\\[4pt]
\large Structural coincidences: $h_\tau\approx\sqrt2\,\alpha$,
$m_\tau/m_e\approx 2\pi/\Omega_\eta$, $m_\tau/m_\mu\approx4\alpha/\Omega_\eta$}
\author{D.~Jaros}
\date{2026-06-14}

\begin{document}
\maketitle
\fi

%──────────────────────────────────────────────────────────────────
\begin{abstract}
We document three structural coincidences between charged lepton masses
and UBT-derived constants. All three involve only $\alpha^{-1}=137.036$
[\Lo, \texttt{alpha\_l1\_proof.tex}] and
$\Omega_\eta=\tfrac{1}{24}-\tfrac{1}{8\pi}$ [\Lz]:
\begin{align*}
  h_\tau &\approx \sqrt{2}\,\alpha & &(-1.0\%), \\
  m_\tau/m_e &\approx 2\pi/\Omega_\eta & &(-3.8\%), \\
  m_\tau/m_\mu &\approx 4\alpha/\Omega_\eta & &(-7.5\%).
\end{align*}
Status: [\MC] --- structural evidence. No derivation from $S[\Theta]$
is available; the EW hierarchy problem ($M_{\rm EW}\ll M_{\rm UBT}$)
must first be resolved. We record these as falsifiable predictions.
\end{abstract}

%──────────────────────────────────────────────────────────────────
\section{UBT Constants}
%──────────────────────────────────────────────────────────────────

\begin{align}
  \alpha^{-1}_{\rm UBT} &= 137.035999177549\ldots &[\Lo]
  \label{eq:alpha}\\
  \Omega_\eta &= \tfrac{1}{24}-\tfrac{1}{8\pi}
  = 0.001877928\ldots &[\Lz]
  \label{eq:omega}\\
  C_Q &= 4 &[\Lo]
  \label{eq:CQ}
\end{align}

%──────────────────────────────────────────────────────────────────
\section{Three Structural Coincidences}
%──────────────────────────────────────────────────────────────────

\begin{observation}[Tau Yukawa $\approx\sqrt2\,\alpha$ — \MC]
The tau-lepton Yukawa coupling $h_\tau=m_\tau\sqrt2/v$ satisfies:
\[
  h_\tau = \frac{m_\tau\sqrt{2}}{v}
  = \frac{1.7769\times\sqrt2}{246.22}
  = 0.010215\ldots
  \approx \sqrt{2}\,\alpha = 0.010320\ldots
\]
Discrepancy: $-1.0\%$. No free parameter used.
\end{observation}

\begin{observation}[Tau-to-electron ratio — \MC]
\[
  \frac{m_\tau}{m_e}
  = \frac{1776.9\,{\rm MeV}}{0.511\,{\rm MeV}}
  = 3477.3
  \approx \frac{2\pi}{\Omega_\eta}
  = \frac{2\pi}{0.001878}
  = 3345.8
\]
Discrepancy: $-3.8\%$.
\end{observation}

\begin{observation}[Tau-to-muon ratio — \MC]
\[
  \frac{m_\tau}{m_\mu}
  = \frac{1776.9\,{\rm MeV}}{105.7\,{\rm MeV}}
  = 16.81
  \approx \frac{4\alpha}{\Omega_\eta}
  = \frac{4/137.036}{0.001878}
  = 15.54
\]
Discrepancy: $-7.5\%$.
\end{observation}

\begin{resultbox}
\textbf{Lepton mass formula (structural, \MC).}
If $m_\tau\approx v\alpha$ and the ratios hold:
\[
  m_\tau : m_\mu : m_e
  \;\approx\;
  v\alpha
  : \frac{v\Omega_\eta}{C_Q}
  : \frac{v\,\Omega_\eta\,\alpha}{2\pi}
  \;=\;
  v\alpha\;:\;\frac{v\Omega_\eta}{4}\;:\;\frac{v\Omega_\eta\alpha}{2\pi}.
\]
All ratios involve only $\alpha$ [\Lo] and $\Omega_\eta$ [\Lz]:
\begin{align*}
  \frac{m_\tau}{m_\mu} &\approx \frac{C_Q\,\alpha}{\Omega_\eta}
  = \frac{4\alpha}{\Omega_\eta}
  = 15.54 \quad(\text{exp.\ }16.81,\; -7.5\%), \\
  \frac{m_\tau}{m_e} &\approx \frac{2\pi}{\Omega_\eta}
  = 3346 \quad(\text{exp.\ }3477,\; -3.8\%).
\end{align*}
\end{resultbox}

%──────────────────────────────────────────────────────────────────
\section{Discussion}
%──────────────────────────────────────────────────────────────────

\begin{remark}[Physical interpretation]
The formula $m_\tau\approx v\alpha$ suggests the tau Yukawa coupling
$h_\tau\approx\sqrt{2}\,\alpha$ is proportional to the EM coupling.
This would arise naturally if the tau mass is generated at one loop
in EM, analogously to the anomalous magnetic moment ($g-2\sim\alpha/\pi$).

The formula $m_e\approx v\,\Omega_\eta\alpha/(2\pi)$ suggests the
electron mass involves two suppressions: one from $\Omega_\eta$
(eta-mode occupation, a $\psi$-sector loop factor) and one from
$\alpha/(2\pi)$ (EM loop factor).
\end{remark}

\begin{openprobbox}
\textbf{What is missing.}

These coincidences are [\MC], not [\Lo]. Three things are needed:

\begin{enumerate}
\item \textbf{VEV hierarchy} $v=246\,{\rm GeV}\ll M_{\rm UBT}$:
  the EW hierarchy problem [\OPEN] must first be resolved.

\item \textbf{Derivation of $h_\tau=\sqrt{2}\,\alpha$} from $S[\Theta]$:
  the tau Yukawa coupling must be shown to equal $\sqrt{2}\,\alpha$
  in the UBT Yukawa sector.

\item \textbf{Derivation of generation structure}:
  why do the three generations correspond to this specific pattern
  $(\alpha, \Omega_\eta/C_Q, \Omega_\eta\alpha/2\pi)$?
\end{enumerate}

Without resolving the hierarchy problem, these remain structural
coincidences without a derivation.
\end{openprobbox}

\begin{remark}[Falsifiability]
The predictions $m_\tau/m_e = 2\pi/\Omega_\eta = 3346$ and
$m_\tau/m_\mu = 4\alpha/\Omega_\eta = 15.5$ are falsifiable.
They are wrong at the 4--8\% level, which may indicate:
(a) higher-order corrections of order $\Omega_\eta$ or $\alpha/\pi$;
(b) the identification is approximate and the true formula differs;
(c) coincidence.
Distinguishing these requires a derivation from $S[\Theta]$.
\end{remark}

\begin{thebibliography}{9}
\bibitem{alpha}
  D.~Jaros, \texttt{alpha\_l1\_proof.tex}, 2026.
\bibitem{layer2}
  D.~Jaros, \texttt{layer2\_kernel\_derivation.tex}, 2026.
\bibitem{higgs}
  D.~Jaros, \texttt{higgs\_yukawa\_status.tex}, 2026.
\end{thebibliography}

\ifdefined\INCLUDEMODE\else
\end{document}
\fi
